\section{Auxiliary markings and the height bound}\label{sec:height}

We prove the nonidentity required by \Cref{thm:interpolation} for the
common relations of \Cref{prop:common-relations}, and then deduce the
height bound.  Single and product monodromy handle the constant relations
and the paired relations away from the finitely many dominating isogeny
correspondences.  The remaining case requires one auxiliary level.

Here is the role of that level.  At a common reduction, Rosati positivity
makes the rational span of the centered and target endomorphism planes
nondegenerate.  We arrange that the sum of their $\Z_\ell$-lattices is
nevertheless saturated of rank four, with degenerate reduction modulo
a fixed prime $\ell$.  A functional identity would force the composite
of relative Frobenius with the inverse of a multiplier-$m$ isogeny
to act by conjugation as one common sign on the rational span.  The positive
sign forces a rational scalar, incompatible with its multiplier $p/m$,
where $p$ is the residue characteristic and $m$ is a fixed
correspondence multiplier prime to $p$.  The negative sign makes the rank-four
lattice an orthogonal eigensummand of a unimodular lattice, whose
reduction modulo $\ell$ must be nondegenerate.  Thus either sign gives
a contradiction.  The marking imposes no restriction on the quaternion
algebra of the target.

\subsection{A finite-field incidence}

The required marking comes from the following incidence statement.
Its acting element must lift to the spin group because symplectic
level markings act on the quadratic system through
$\Sp_4\longrightarrow\SO_5$.

\begin{lemma}\label{lem:finite-field-incidence}
Let $k$ be a finite field of odd characteristic, let $(V,q)$ be a nondegenerate
quadratic space of dimension five, let $A\subset V$ be a nondegenerate
$2$-plane, and let $T\subset V$ be any $2$-plane.  There is an element
\[
 u\in\operatorname{im}\bigl(\Spin(V)(k)\longrightarrow\SO(V)(k)\bigr)
\]
such that $A\cap uT=0$ and $A+uT$ is a degenerate $4$-plane, with radical
of dimension one.  The conclusion includes restricted quadratic forms on
$T$ of ranks zero, one, and two.
\end{lemma}

\begin{proof}
Write $b(v,w)=q(v+w)-q(v)-q(w)$.  The nondegenerate ternary space
$A^\perp$ is isotropic over $k$, so it has a decomposition
\[
 A^\perp=\langle e,f\rangle\perp\langle c\rangle,
 \qquad q(e)=q(f)=0,\quad b(e,f)=1,\quad q(c)=\delta\ne0.
\]
Here and below angle brackets around vectors mean their linear span.  For
any $v,w\in A$, the plane
\[
 T_* = \langle e+v,c+w\rangle
\]
is disjoint from $A$, and $A+T_*=e^\perp$.  This hyperplane has radical $ke$.
We can arrange that $T_*$ and $T$ have isometric restricted forms.  A
nondegenerate binary form over $k$ represents every nonzero scalar: in the
split case use the form $xy$, and in the anisotropic case use surjectivity
of the norm $k_2^\times\to k^\times$.  The value zero is represented by the
zero vector.  Consequently the following choices are possible:
\begin{enumerate}[label=(\roman*),leftmargin=*]
\item If $q|_T=0$, take $v=0$ and $q(w)=-\delta$.
\item If $q|_T$ has rank one, write it as $\operatorname{diag}(0,\beta)$
with $\beta\ne0$, take $v=0$, and choose $q(w)=\beta-\delta$.
\item If $q|_T$ has rank two, choose a vector of norm $\delta$ in $T$ and
an orthogonal complement of norm $\alpha\ne0$.  Take $q(v)=\alpha$ and
$w=0$.  Then $q|_{T_*}=\operatorname{diag}(\alpha,\delta)$.
\end{enumerate}
Witt extension, which applies also to isometries between degenerate
subspaces of a nondegenerate space in characteristic different from two
\cite[Theorem~8.3]{EKM2008}, extends the chosen isometry $T\to T_*$ to an element of $\operatorname{O}(V)(k)$.
If its determinant is $-1$, compose it with reflection in an anisotropic
vector of $T_*^\perp$.  Such a vector exists: a three-dimensional subspace
of a nondegenerate five-dimensional space cannot be totally isotropic.
The reflection fixes $T_*$ pointwise.  We have therefore obtained
$u_0\in\SO(V)(k)$ with $u_0T=T_*$.

It remains to correct the spinor norm.  The setwise stabilizer of every
$2$-plane $E$ in $\SO(V)(k)$ has surjective spinor norm.  To see this,
choose a hyperbolic plane $\langle r,s\rangle$, normalized by $b(r,s)=1$,
and consider
\[
 h_z(r)=zr,\qquad h_z(s)=z^{-1}s,\qquad
 h_z|_{\langle r,s\rangle^\perp}=1\quad(z\in k^\times).
\]
If $R_v$ denotes reflection in $v$, then
$h_z=R_{r+s}R_{r+zs}$, so its spinor norm is $z$ modulo squares.
When $E$ is nondegenerate, take this hyperbolic plane inside the
nondegenerate ternary space $E^\perp$; then $h_z$ fixes $E$ pointwise.
When $E$ is degenerate, write $E=\langle r,a\rangle$ with
$0\ne r\in\operatorname{rad}(E)$.  Choose $s_0$ with
$b(r,s_0)=1$ and $b(a,s_0)=0$, and put $s=s_0-q(s_0)r$.
Then $q(s)=0$ and $a\perp\langle r,s\rangle$, so $h_z$ preserves
$E$ setwise.  This also treats the totally isotropic case.

Apply the preceding construction to $E=T_*$, choosing the spinor norm
of $h_z$ to be the inverse of that of $u_0$.  The element $u=h_zu_0$ still
maps $T$ to $T_*$ and has trivial spinor norm.  To see the lifting
assertion directly,
write an orthogonal transformation of determinant one as an even product
of reflections \cite[Theorem~1.32]{SchulzePillot2020}.  The corresponding product in the Clifford algebra has
Clifford norm equal to the product of their quadratic norms.  If this
product is a square $c^2$, rescaling the Clifford element by $c^{-1}$
gives a norm-one element, hence an element of $\Spin(V)(k)$ with the
same orthogonal action \cite[Corollary~9.6 and Definition~9.7]{SchulzePillot2020}.
Conversely a spin lift has trivial spinor norm.
Thus $u$ has the required lift.  In dimension five over $k$, this spin
group is the split group $\Sp_4$.
\end{proof}

\subsection{Fixing the level and the exceptional primes}

We first specify the order of the fixed choices.  Apply
\Cref{lem:finite-covers} to obtain a smooth connected fine-level curve
$S_*$ and its family.  If the quaternionic-division locus is finite, both
the finiteness
assertion and the height bound below hold after choosing constants.  We
therefore assume in the construction that the locus is infinite, and
choose $t_{0,*}$ away from the finite omissions.  Fix its labeled
endomorphism pair and all the conjugates of
these data.  Apply \Cref{prop:finite-correspondences} to the finite list of
underlying image curves.  Choose integral representatives of the resulting
rational similitudes and let $\mathcal M$ be their finite set of positive
polarized multipliers.  This step precedes the auxiliary level: passing to
that level does not enlarge the list of multipliers required for the
underlying image curves.

For a rational endomorphism plane $E_t$, let
\[
 L_t=E_t\cap\Lambda_t
\]
be its saturated lattice in the integral quadratic local system of
\Cref{lem:qm-plane}.  The integral lattice is understood with the fixed
finite set of bad denominator primes removed.  Saturation implies that
$L_t/\ell L_t$ is a $2$-plane in $\Lambda_t/\ell\Lambda_t$ at every remaining
prime $\ell$.  This holds even when its quadratic form becomes degenerate.
The small generating pair from \Cref{prop:small-vectors} need not be a basis
of $L_t$ and will not be used to define its reduction.

\begin{proposition}\label{prop:marking}
There is a fixed odd prime $\ell$ and a fixed geometrically connected finite
\'{e}tale cover $S\to S_*$ parameterizing symplectic markings at level $\ell$,
with a fixed lift $t_0$ of $t_{0,*}$, such that every quaternionic-division
point of $S_*$ outside a finite set has a lift $t\in S$ satisfying the
following condition.  In the common quadratic space supplied by the
markings,
\begin{equation}\label{eq:marked-incidence}
 \dim\bigl(\overline L_{t_0}+\overline L_t\bigr)=4,
 \qquad
 q|_{\overline L_{t_0}+\overline L_t}\text{ is degenerate}.
\end{equation}
The prime $\ell$ is prime to every $m\in\mathcal M$.  At every prime
$p\ne\ell$ of good reduction for these fixed models where $t$ and $t_0$
have the same marked reduction, the corresponding two saturated
$\Z_\ell$-lattices have sum $M$ satisfying
\begin{equation}\label{eq:primitive-four-lattice}
 M=\Lambda_\ell\cap(M\otimes_{\Z_\ell}\Q_\ell),
 \qquad \rank M=4,
 \qquad M/\ell M\text{ is degenerate}.
\end{equation}
These assertions are preserved by simultaneous field conjugation.
\end{proposition}

\begin{proof}
The geometric monodromy subgroup in $\Sp_4(\Z)$ is finitely generated,
because $S_*^{\mathrm{an}}$ has finitely generated fundamental group, and
is Zariski dense by \Cref{prop:monodromy}.  Strong approximation for a
finitely generated Zariski-dense subgroup of a simply connected absolutely
almost simple group gives surjectivity onto $\Sp_4(\mathbb F_\ell)$ for all
sufficiently large primes $\ell$ \cite[Theorem, p.~515]{MVW1984}.  Choose
such an $\ell$, avoiding $2$, the discriminant of the ambient quadratic
lattice, the discriminant of $L_{t_{0,*}}$, and every multiplier in
$\mathcal M$.  After a fixed extension containing the relevant root of
unity, full symplectic markings with the prescribed pairing form a torsor
under $\Sp_4(\mathbb F_\ell)$.  Surjective geometric monodromy says exactly
that its pullback to $S_*$ is geometrically connected.

Fix a marking above $t_{0,*}$.  Its plane is nondegenerate modulo $\ell$.
For any target, its saturated plane still has dimension two modulo
$\ell$, and \Cref{lem:finite-field-incidence} supplies a change of marking
in the image of $\Sp_4(\mathbb F_\ell)$ giving
\eqref{eq:marked-incidence}.  All markings occur on the same cover $S$.
Deleting finitely many points to choose frames and affine charts excludes
only finitely many points downstairs.  In particular, $\ell$ has not been
chosen to avoid the discriminants of the varying target planes.

At a common marked reduction, smooth proper base change identifies each
integral $\ell$-adic cohomology lattice with that of the common reduction.
The mod-$\ell$ identifications agree with the markings.  Take
$\Z_\ell$-bases of the two saturated planes.  Their four reductions are
independent by \eqref{eq:marked-incidence}, so some $4$-row minor of their
coordinate matrix is an $\ell$-adic unit.  These four vectors extend to a
basis of $\Lambda_\ell$; their sum $M$ is a direct summand, and therefore
satisfies the saturation equality in \eqref{eq:primitive-four-lattice}.
Its reduction is precisely the degenerate space in
\eqref{eq:marked-incidence}.  This proves local saturation at $\ell$;
no global saturation assertion about the sum of integral lattices is
needed.  Simultaneous conjugation carries both markings, both lattices,
and their common pairing to the conjugate data, preserving independence
and degeneracy.
\end{proof}

We have now fixed the incidence condition and its integral consequence.
Enlarge the fixed field $K$ for the cover, the chosen lift $t_0$, and its
labeled pair, and use the notation of \Cref{sec:arithmetic} on $S$.
Before grouping proximity places, we make one further finite exclusion:
every isogeny of an allowed multiplier between the reduced centers must
lift to one of finitely many characteristic-zero isogenies between them.  The next
lemma provides this assertion for the maps themselves.

\begin{lemma}\label{lem:fixed-isogeny-reduction}
Fix finitely many pairs of principally polarized abelian surfaces over a
number field and finitely many positive integers $m$.  Outside a fixed
finite set of rational primes, every polarized isogeny of multiplier $m$
between the geometric reductions of a fixed pair is the specialization of
a characteristic-zero polarized isogeny of that multiplier between the
pair.  The set of these characteristic-zero maps is finite.  The assertion
concerns the maps themselves, and holds for supersingular reductions as
well.
\end{lemma}

\begin{proof}
It suffices to treat one pair $(A_0,\lambda_0)$, $(A'_0,\lambda'_0)$ and
one $m$.  Choose models over a localization of the ring of integers, with
good reduction and with $m$ invertible.  A polarized isogeny
$\varphi$ with $\varphi^*\lambda'_0=m\lambda_0$ has degree $m^2$ and kernel
contained in $A_0[m]$.  After a fixed finite extension and localization,
all subgroup schemes of this fixed order and the required isotropic type
in $A_0[m]$ are the spreads of the finitely many characteristic-zero
ones.  Indeed $A_0[m]$ is finite \'{e}tale, and after splitting it these
are subgroups of one fixed finite group.

Each possible kernel gives a fixed polarized quotient
$(D,\lambda_D)$, with the quotient map pulling $\lambda_D$ back to
$m\lambda_0$.  The remaining data are polarized isomorphisms
$(D,\lambda_D)\simeq(A'_0,\lambda'_0)$.  Here isomorphisms are group isomorphisms of abelian schemes, preserving
the zero section and the polarizations.  Their scheme is unramified by
rigidity and is of finite type: for relatively ample line bundles
representing the polarizations, the restriction of their product to
such a graph is numerically twice the source line bundle, and hence
has a fixed Hilbert polynomial.  Thus these graphs lie in one relative
Hilbert scheme.  The fibers of the polarized-isomorphism scheme are
torsors under the finite polarized
group-automorphism groups \cite[Proposition~17.5]{Milne1986}, so the
scheme is quasi-finite.

A quasi-finite scheme of
finite type over a ring of integers,
with finite generic fiber, becomes the spread of that generic fiber
after deleting finitely many primes.  For completeness, discard its
finitely many vertical irreducible components; its finitely many generic
points then spread to a finite scheme after localization, and deleting
the support of the remaining complement gives the assertion.  Further
localization makes this scheme finite \'{e}tale.  After a fixed extension splitting its generic fiber, this finite
\'{e}tale scheme is a disjoint union of sections: each component has
generic degree one over the normal base.  Its geometric points are
therefore exactly the specializations of the fixed characteristic-zero
maps, including every polarized-automorphism choice.  Applying this to the finite list of kernels
proves the lemma.
\end{proof}

Include the primes excluded by \Cref{lem:fixed-isogeny-reduction}, those
dividing the elements of $\mathcal M$, and $\ell$ in the fixed exceptional
set of \Cref{prop:common-relations}.  Here the center pairs are the finitely
many pairs of conjugates of $t_0$.  The possible center isogenies are now
fixed algebraic data.  Take the full finite list of these maps between
every ordered pair of absolute conjugates of the center; the list is
stable under simultaneous absolute Galois conjugation.  Enlarging $K$
once more defines all of them.  This enlargement adds no geometric
center to the list, since the absolute conjugates were included already.
The local Frobenius lift in \Cref{prop:common-relations} need not fix
$K$: it permutes the conjugate centers and maps in this fixed list.
At an embedding $\rho$ of the labeled coefficient field, a selected
specialized center map is $\rho(\varphi_0)$ for a member $\varphi_0$
of that list.  Subdividing by this member makes the algebraic choice
common to the selected embeddings.  Upon any further coefficient-field
extension we retain all extensions of each selected embedding, so
the map and its specialization remain compatible and the total weights
are unchanged.  A
chosen multiplier-$m$ isogeny between target fibers also requires only a
bounded relative field extension.  Its kernel is one of a uniformly
bounded number of subgroups of the fixed-size group $A_t[m]$, and the
polarized identifications form a torsor under a finite automorphism group
of bounded order in dimension two.  For the latter bound, such a group
acts faithfully on the rank-four integral Betti lattice, and finite
subgroups of $\operatorname{GL}_4(\Z)$ have bounded order.  Thus these
choices preserve all
polynomial degree bounds in \Cref{sec:arithmetic}.

\subsection{Nonidentity on every branch}

We use three lists of equations.  Constant relations will be excluded
by single monodromy, and paired relations outside the finite
correspondence locus by product monodromy.  On that locus we first
convert the paired relation to one involving a single transport
matrix; the marking will then exclude an identity.

The notation $B=(B_1,B_2)$ and $P=(P_1,P_2)$ means the labeled centered
and target pairs, and $Y(a)$ transports back to the center.  We retain
all four entries of each matrix equation: in the exceptional case they
will force the same sign on both entire planes.  In every case include
the algebraic chart equations from
\Cref{prop:common-relations}: if $z$ denotes affine coordinates on $S$ and
$z=\zeta(a)$ is the actual inverse germ of the chosen parameter, include
$z-\zeta(a)=0$.  For a second argument include
$z'-\zeta'(aR)=0$.  The ambient algebraic conditions put $z,z'$ on the
respective curves and impose $a=x(z)$ and $aR=x'(z')$ when appropriate.
Only fixed denominators needed for these affine charts are cleared.

The first list, used at archimedean or fixed exceptional places, is
\begin{equation}\label{eq:constant-system}
 \pair{B}{Y(a)P}=c,
\end{equation}
where the $2$-by-$2$ rational matrix $c$ is common to the chosen group.
The second list is
\begin{equation}\label{eq:paired-system}
 \pair{B}{Y(a)P}=\pair{B'}{Y'(aR)P'}.
\end{equation}
It is used when the target pair of image points is outside all the
finitely many dominating isogeny correspondences of
\Cref{prop:finite-correspondences}.

For the third list the target pair belongs to one of those
correspondences.  We transport its primed vectors to the unprimed fiber
using an isogeny, and do the same at the centers.  To compare these two
operations, their isogenies must have the same specialization.
Choose a polarized isogeny $\varphi_s$ of multiplier
$m\in\mathcal M$ between the targets, and write $I_s$ for its normalized
isometry on the quadratic systems.  First adjoin its field of definition
to the paired target field and extend every already selected embedding
by all its extensions.  This costs only a bounded relative degree and
preserves every weighted local sum.  The previously chosen Frobenius
identities restrict to the original labeled tuples, so remain true
under these extensions without a new choice of Frobenius lift.
At a place of this expanded group the unprimed center and target reduce to the
same marked point $r$ of the good fine model, and hence to the same
universal fiber $A_r$.  The chosen local automorphism has residue action
$\operatorname{Frob}_p$, so the primed center and target both reduce to
$r^{(p)}$ and to $A_r^{(p)}$.  These are the conjugate identifications
of the same source fiber, not independently chosen isomorphisms.
The target isogeny extends between the good abelian models; its
specialization is an actual multiplier-$m$ map
$A_r\to A_r^{(p)}$.  By \Cref{lem:fixed-isogeny-reduction} this map
is the specialization of one fixed center isogeny $\varphi_0$, with
normalized isometry $I_0$.  Only now subdivide the expanded embedding
set by this finite choice from the conjugate-stable list of actual
maps, including each map's polarized quotient identification.

The center and target comparison maps now identify both isogenies with
the same specialized map.  Functoriality in
\Cref{prop:good-comparison} therefore gives, at this target,
\begin{equation}\label{eq:isogeny-transport}
 I_0Y(a)=Y'(aR)I_s.
\end{equation}
Put $D_i=I_0^{-1}B'_i$ and $\widehat P_j=I_s^{-1}P'_j$ at the target.
The third list is the single-argument system
\begin{equation}\label{eq:exceptional-system}
 \pair{B}{Y(a)P}=\pair{D}{Y(a)\widehat P}.
\end{equation}
In this system $D$ is fixed and $\widehat P$ is an additional pair of
vector coordinates in the unprimed frame; the second argument can be
dropped.  Its target value is the realization of actual rational
endomorphisms, with denominators bounded in terms of $m$ and with the same
Rosati norm bounds as $P'$.  More explicitly, if
$\widehat\alpha=\varphi_s^{-1}\alpha'\varphi_s$, then
$m\widehat\alpha=\varphi_s^\dagger\alpha'\varphi_s$ is integral and
$q_{\lambda_s}(\widehat\alpha)=q_{\lambda'_s}(\alpha')$.
Thus \Cref{prop:dr-size,prop:common-relations}
apply to this new tuple.  No isogeny is assumed to exist along a containing
algebraic variety used in interpolation.

\begin{proposition}\label{prop:nonidentity}
For a selected marked target as in \Cref{prop:marking}, each applicable
list above, formed from a nonempty selected group of places, satisfies
the all-branch nonidentity hypothesis of
\Cref{thm:interpolation}.  More precisely, let $W$ be any irreducible
algebraic subvariety of the indicated affine ambient conditions which
contains the target and on which $a$ is nonconstant.  On every branch
over the generic point of every component of $W\cap\{a=0\}$, at least
one equation in the applicable list, including the chart equations,
is not an identity.
\end{proposition}

\begin{proof}
\emph{From a branch identity to monodromy.}
Suppose that every equation is an identity on such a
branch.  On the normalization near a general point of its divisor, $a$
vanishes and all affine coordinates, including $R$ when present, are
regular.  The series in $a$ and $aR$ therefore converge after a complex
embedding in a neighborhood of that point.  In a transverse parameter on the normalization, the formal identity
makes every coefficient zero at the generic point of the divisor.
Convergence therefore gives an identity on a nonempty analytic open.
Any fixed denominator cleared in a chart is nonzero at the target,
hence is a nonzero function on $W$ and has finite order on each
normalization branch.  It cannot turn a nonzero series into an identity.  If one of the chart
equations is not an identity we already have the desired conclusion;
otherwise they identify the series arguments with the actual
projections to $S$ and $S'$.  The first projection is nonconstant since
$a$ is nonconstant.

Pass to a smooth dense open of $W$ on which these projections and the
algebraic frames are defined.  An irreducible complex algebraic variety
has connected smooth locus, so the analytic identity continues on its
universal cover.  At a fixed point of this open the vector coordinates
are single-valued, whereas continuation changes the transport matrices
by the monodromy of the pulled-back systems.  A dominant map to a curve
has fundamental-group image of finite index after deleting suitable
closed subsets and passing to a resolution; thus
\Cref{prop:monodromy} applies on this open.  The vector pairs are
generically independent, because they are independent at the target.
A pair minor nonzero at the target is not identically zero on $W$,
so its nonvanishing locus meets this analytic open even if the pair
vanishes on the divisor itself.
They need not remain independent at every boundary point of $W$.

\emph{Constant relations and independent projections.}
For \eqref{eq:constant-system}, fix a general point of this open.  The
orbit of the transport matrix under continuation has Zariski closure a
translate of $\SO_5$.  For nonzero vectors $b,v$ in its standard
representation, the function $g\mapsto\pair{b}{gv}$ is not constant.
Indeed, the span of the differences $gv-v$ is invariant, and is the whole
representation by irreducibility and the absence of invariant vectors.
A constant matrix coefficient would annihilate this span.  Applying this
to $B_1$ and $P_1$ contradicts the continued identity.

For \eqref{eq:paired-system}, if the second projection is constant, only
the first system varies and the same argument applies.  If it is
nonconstant, \Cref{prop:monodromy} gives product monodromy unless the
closure of the paired image is contained in a dominating isogeny
correspondence.  In the latter case
\Cref{prop:finite-correspondences} puts the image in the fixed finite list;
since these correspondences are closed, the target also belongs to that
list, contrary to this case's defining assumption.  For product
monodromy, fix the second transport matrix and vary the first.  The
resulting constant matrix coefficient is again impossible.

\emph{The exceptional relation forces a common sign.}
Consider now \eqref{eq:exceptional-system}.  At a general point $z$ of
$W$, write $Y_z:\Vcal_z\to\Vcal_{t_0}$ for transport.  The continued
identity holds with transport $gY_z$ for every
$g\in\SO(\Vcal_{t_0})$.  The standard representation is absolutely
irreducible, so Burnside's theorem says that these matrices span
$\Hom(\Vcal_z,\Vcal_{t_0})$.  For each $i,j$, the linear functional
\[
 M\longmapsto\pair{B_i}{MP_j}-\pair{D_i}{M\widehat P_j}
\]
therefore vanishes on every linear map $M$ between these two spaces.
The corresponding rank-one tensors are equal.  Identifying their first
factors with covectors by the fixed center trace pairing, we obtain
\begin{equation}\label{eq:burnside-tensors}
 B_i\otimes P_j=D_i\otimes\widehat P_j
 \qquad(1\le i,j\le2).
\end{equation}
These are polynomial equations in the algebraic vector coordinates.
They hold on a dense open of $W$, and therefore at the target itself.
There, all four pairs of factors are nonzero.  The $(1,1)$ equation
gives $D_1=c_0B_1$ and $\widehat P_1=c_0^{-1}P_1$ for some nonzero
scalar $c_0$; the $(2,1)$ and $(1,2)$ equations give the same scalar for
the other indices.  The self-pairing of $B_1$ is a nonzero rational
number, and is unchanged by conjugation or the isometry $I_0$.  Thus
$c_0^2=1$, and
\begin{equation}\label{eq:common-sign}
 D_i=\eps B_i,\qquad \widehat P_j=\eps P_j,
 \qquad \eps\in\{1,-1\}.
\end{equation}
The $(2,1)$ and $(1,2)$ equations are essential here: they couple the
two labels and give a single sign for both pairs.

\emph{Specialization to a Frobenius quasi-isogeny.}
We now use that the vectors in \eqref{eq:common-sign} come from
endomorphisms.  De Rham realization on the rational endomorphism algebra
of a characteristic-zero abelian variety is faithful: base change to
$\C$, compare with Betti cohomology, and use the faithful action on its
rational first homology.  Conjugation by an isogeny preserves the Rosati
self-adjoint subspace.  Consequently \eqref{eq:common-sign} asserts
equalities of actual rational endomorphisms in the center and target
fibers.  These equalities specialize under good reduction.

Choose any place in the nonempty selected group and let $p$ be its
residue characteristic.  The unprimed center and target have a common
marked reduction $\overline A$, and the primed reductions identify with
its $p$-twist.  Write $F:\overline A\to\overline A^{(p)}$ for relative
Frobenius, and write $\overline\varphi_0$ for the common specialization of
$\varphi_0$ and $\varphi_s$.  The map
\[
 u=\overline\varphi_0^{-1}F\in\End^0(\overline A)
\]
is a self-quasi-isogeny with polarized multiplier $p/m$.  Relative
Frobenius transports the unprimed labeled endomorphisms to their primed
specializations.  The specialized equalities \eqref{eq:common-sign}
therefore say that the conjugation action $g$ of $u$ on the quadratic
$\Q_\ell$-space acts by $\eps$ on both specialized endomorphism planes.
The sign action extends to their saturated $\Z_\ell$-lattices; the
chosen small pairs need not be bases of those lattices.
By \Cref{lem:qm-plane}, $g$ belongs to $\SO_5$.

\emph{Excluding the two signs.}
Let $T$ be the rational span of those two planes in
$\End^0(\overline A)$.  Its $\ell$-adic realization has dimension four by
\Cref{prop:marking}.  The trace form on its rational symmetric
endomorphisms is positive definite by Rosati positivity in arbitrary
characteristic \cite[Theorem~17.3]{Milne1986}; hence $T$ is nondegenerate
over $\Q$, and its realization
$T_\ell$ is nondegenerate over $\Q_\ell$.  The one-dimensional perpendicular
is consequently $g$-stable.  Since $\det g=1$ and
$\det(g|_{T_\ell})=\eps^4=1$, $g$ acts as $+1$ on that perpendicular.

If $\eps=+1$, it follows that $g=1$.  The kernel of the conjugation
representation of $\GSp_4$ on this quadratic space is its scalar center
by \Cref{lem:qm-plane}.  Put $\mu=p/m$.  With the covariant convention
of \Cref{sec:geometry}, the action on first cohomology is
\[
 T_u=\mu(u^*)^{-1}=(u^\dagger)^*;
\]
the equality follows from $u^\dagger u=[\mu]$.  Thus
$(u^\dagger)^*=c\id$ for some $c\in\Q_\ell$.  The ordinary
cohomological trace of the rational quasi-endomorphism $u^\dagger$
is rational \cite[Proposition~12.9]{Milne1986}, so
$4c=\Tr(u^\dagger)\in\Q$.  Faithfulness gives
$u^\dagger=c\id$, and applying Rosati gives $u=c\id$.
Its multiplier would satisfy $c^2=p/m$.  This is impossible because
$p\nmid m$, so the $p$-adic valuation of $p/m$ is odd.

If $\eps=-1$, $g$ is an involution with negative eigenspace exactly
$T_\ell$.  Both $F$ and $\overline\varphi_0$ have multipliers prime to
$\ell$, so their actions are invertible on the integral Tate modules;
conjugation by $u$ preserves the integral quadratic lattice
$\Lambda_\ell$.  This lattice is unimodular by the choice of $\ell$.
Since $\ell$ is odd, the integral projectors $(1+g)/2$ and $(1-g)/2$
give an orthogonal direct sum
\[
 \Lambda_\ell=\Lambda_+\perp\Lambda_-,
 \qquad \Lambda_-=\Lambda_\ell\cap T_\ell.
\]
Each summand is unimodular, so $\Lambda_-/\ell\Lambda_-$ is nondegenerate.
But \Cref{prop:marking} identifies $\Lambda_-$ with the sum $M$ of the
two marked saturated plane lattices, whose reduction is degenerate by
\eqref{eq:primitive-four-lattice}.  This is the required contradiction.

Both cases contradict the assumed identity.  The use of Rosati positivity
and rational endomorphism trace is valid for every reduction type; no
classification or ordinarity assertion about Frobenius has been used.
\end{proof}

\subsection{Application of interpolation and descent}

\begin{theorem}\label{thm:height}
Fix a reduced irreducible closed Hodge-generic curve $C\subset\Acal_2$
over $\Qbar$, and let $C^\nu$ be its normalization.  For a fixed
nonnegative ample logarithmic height $h$ on the smooth projective
completion of $C^\nu$ and a fixed field of definition $K$, there are
constants $C_1,C_2>0$ such that every point $s\in C^\nu(\Qbar)$ whose
coarse image has full geometric rational endomorphism algebra an
indefinite quaternion division algebra over $\Q$ satisfies
\[
 \max\{2,h(s)\}\le C_1\max\{2,[K(s):K]\}^{C_2}.
\]
The same assertion holds on the fixed fine-level curves used above,
and for a fixed projective height of the coarse points.  The constants
do not depend on the quaternion algebra, its discriminant, or
the target endomorphism order.
\end{theorem}

\begin{proof}
If the target locus is finite, enlarge the constants to cover it.
Otherwise use the initial family and center, the finite multiplier list,
and then the fixed prime and connected marking cover, in the order already
specified.  For each target choose one lift satisfying
\Cref{prop:marking}; this does not enlarge any fixed datum.  Remove the
finite exceptional targets of the affine charts, and suppose that its
height $h$ exceeds a sufficiently large fixed polynomial in
$d=\max\{2,[K(t):K]\}$.  The fields generated by the labeled single or paired point tuples in
\Cref{prop:common-relations} have degree polynomial in $d$.
A normal closure is used only to index selected embeddings.  A chosen
fixed-multiplier isogeny requires only a bounded relative extension over
the field of its two point arguments, so adjoining it preserves this
polynomial bound.  Selected embeddings are extended by all extensions;
the new weights above an old embedding sum to its old weight, preserving
the normalized smallness sum.

\emph{Arithmetic inputs.}
Apply \Cref{prop:common-relations}, with the fixed exceptional primes
chosen above.  Partition first into its constant and Frobenius groups,
and, if needed, subdivide the latter by the finite center-isogeny choice.
There are polynomially many groups, and one supplies a common list of
equations and a selected weighted set $\mathcal V$ with
\[
 \sum_{v\in\mathcal V}w_v\log\frac1{\abs{a}_v}
       \ge \frac{h}{Q_1(d)},\qquad
 \sum_{v\in\mathcal V}w_v\log^+\norm{U}_v
       \le Q_2(d)(1+\log^c h).
\]
The tuple $(a,U)$ has global logarithmic height at most $Q_3(d)h$.
The coefficient degrees are $O(j+1)$, and the complete affine
coefficient tuple through order $l$, including scale, has logarithmic
height at most $Q_4(l,d)(1+\log^c h)$.  The tails satisfy the hypotheses
of \Cref{thm:interpolation}.  These assertions also hold
for the transformed target pair $\widehat P$ in the exceptional case:
its fixed denominator and preserved Rosati norms give exactly the size
hypotheses of \Cref{prop:dr-size}, and only the original ordinary
connection in $a$ is used.  In the other paired case $\abs{R}_v=1$ at
every selected place.  Thus its two ordinary systems are evaluated at
$a$ and $aR$, as required.  The chart equations use fixed algebraic
ordinary germs.

\emph{Nonidentity and interpolation.}
Start the interpolation argument inside the fixed affine algebraic
conditions on the base coordinates described above.  They have bounded
complexity.  No inverse Gram matrices or inverse minors of the target
endomorphism pairs have been adjoined.  Every irreducible subvariety
encountered in the dimension descent contains the target.  If $a$ varies
there, \Cref{prop:nonidentity} supplies nonidentity on every relevant
branch, including branches on which a chart equation already fails.
If $a$ is constant, the stopping case in
\Cref{thm:interpolation} applies directly.

\Cref{thm:interpolation} now gives $h\le C_3d^{C_4}$ for the selected
lifts.  Its constants are uniform because the cover, centers, allowed
ordinary systems, and isogeny multipliers are fixed.

\emph{Descent from the selected lifts.}
The degree of the marking cover is fixed.  If $t_*$ is the image of $t$,
then
\[
 [K(t):K]\le (\deg(S\to S_*))[K(t_*):K].
\]
Choose the ample height upstairs as the height of the pullback of an
ample divisor downstairs; it equals the downstairs height up to a
bounded error, and arbitrary fixed ample heights on these curves are
comparable up to fixed positive constants and bounded errors.  Thus the
bound for the selected lift descends to $t_*$.  The further normalization
and fine-level maps descend in the same way by \Cref{lem:finite-covers}.
For any other lift $\widetilde t$ in a fixed cover, height comparison
and $[K(t_*):K]\le[K(\widetilde t):K]$ give the same numerical bound.
Thus only the selected lifts require \eqref{eq:marked-incidence}.
At singular coarse points there are finitely many normalization
preimages, all included among the finite exceptions.  Increasing the
constants includes all finitely many omitted points.
The absolute height bound for a selected lift also holds for all its
conjugates over the fixed data field, even when those conjugates were not
selected relative to the original marking of $t_0$; no fresh incidence
choice is needed to deduce that numerical bound.  All fixed enlargements
of the data field $K$ change the degree bounds by at most fixed factors.
\end{proof}
